\section{A delayed quadratic twist on the separated residual locus}
\label{sec:cubic-delayed-quadratic-twist}

We use the unitary normalization of AC-G78--AC-G80.  Thus
\(\epsilon(\operatorname{Frob}_l)=l^{-1}\), the normalized
two-dimensional parameter has determinant \(\eta_K\), and the
normalized three-dimensional parameter is \(\epsilon r_\sigma\).
All coefficient fields may be enlarged finitely.  The two quadratic
extensions below are chosen before the nonabelian Chebotarev prime
\(q\); no bound for that prime is obtained here.

\begin{theorem}[AC-G81: removing the scalar-witness obstruction after a delayed quadratic base change]
\label{thm:G81}
Put \(p=269\), \(K=\mathbf Q(\sqrt{-3})\), and let \(\chi\) be the
fixed anticyclotomic cubic character of conductor \(269\).
Let \(e\ge1\).  Retain the actual ordinary modular endpoint input
of AC-G78: an ordinary cuspidal eigenform \(f_0\) of weight
\(2e+1\), nebentype \(\eta_K\), with underlying newform level
prime to \(p\), and its normalized representation
\[
 V=V(f_0)\otimes\tau,\qquad
 \det V=\eta_K,\qquad
 \tau\in\{1,\text{the quadratic character of conductor }p\}.
\]
Assume \(V|_{G_K}\) is absolutely irreducible and
\[
 \overline V^{\rm ss}
   =(\eta_K\overline\gamma)\oplus\overline\gamma^{-1},
 \qquad \gamma=\psi\epsilon^e,\qquad
 \overline\gamma^{\,2}|_{G_K}\ne1,                 \tag{G81.1}
\]
with the ordinary local shape stipulated in AC-G78 and the actual
coefficient-prime conditions used in AC-G79.
In particular the theorem applies to inherited actual Type II
endpoints whenever AC-G68, the good-twist construction and
AC-K29-nonCM supply these inputs and (G81.1) holds.

There are two fresh imaginary quadratic fields
\(D=\mathbf Q(\sqrt{-r_0})\) and
\(K'=\mathbf Q(\sqrt{-r_1})\), and a triquadratic CM field
\[
 F=KK',\qquad L=FD,\qquad L^+=L\cap\mathbf R,       \tag{G81.2}
\]
with the following properties.
The extension \(L/L^+\) is unramified at every finite place,
\([L^+:\mathbf Q]=4\), and \(L=L^+K'\).
Both \(D\) and \(K'\) split at \(p\).
There is an ordinary automorphic representation \(\sigma_L\)
of the definite rank-three unitary group for \(L/L^+\), of the
original parallel symmetric weight
\(\mu=(e-1,0,1-e)\), such that the normalized parameter
\(R_L=\epsilon r_{\sigma_L}\) is absolutely irreducible and
\[
 \overline R_L^{\rm ss}
  =(\overline\chi^{\,j}\overline\gamma)|_{G_L}
       \oplus1\oplus
    (\overline\chi^{\,j}\overline\gamma^{-1})|_{G_L},
       \qquad j\in\{1,2\}.                        \tag{G81.3}
\]
At each coefficient-prime place it has the exact ordered flag
of AC-G79, with graded characters
\[
 u_1\chi^j\tau\epsilon^e,\qquad
 u_2,\qquad
 u_3\chi^j\tau\epsilon^{-e},                       \tag{G81.4}
\]
where the \(u_i\) are unramified.  Its monodromy there is zero.
At the auxiliary \(q\)-places its inertia is finite and
monodromy is zero.

The construction removes the sign-witness obstruction on the
separated locus (G81.1), including even \(e\).  It does not assert
that an auxiliary twist of \(\psi\) satisfies the old
\(\mathrm{GD.par}\) or defines an old deformation datum.
It also does not cover branches on which
\(\overline\gamma^{\,2}|_{G_K}=1\).  In particular, for principal
\(\psi\), the residues \(e\equiv0,134\pmod{268}\) are not covered
by this argument.

There is no assertion of a bound for \(q\), original-level
descent, per-curve \(J\)-containment, or preservation of the
numerical \(\mathrm{Frame.Ok}\), \(\mathrm{EC}\), \(\mathrm{Big}\)
and counting-budget inequalities after adjoining \(D\).
\end{theorem}

\begin{proof}
\emph{The fresh twist and its literal scalar witness.}
Choose a finite set \(S\) containing \(2,3,p,269\), the primes
ramified in \(V,\psi,\chi,\tau\), and all other fixed old data.
Dirichlet's theorem supplies a prime \(r_0\notin S\) with
\[
                         r_0\equiv-1\pmod{4p}.
                                                        \tag{G81.5}
\]
Let \(\delta=\eta_{-r_0}\) be the quadratic character of \(D\).
It is odd, ramifies only at \(r_0\), and is trivial at \(p\):
its local character is unramified and its Frobenius value is
\(({-r_0}/p)=1\).  Thus \(D\) splits at \(p\).

Twist the actual cuspidal representation of \(f_0\) by
\(\delta\circ\det\).  On automorphic functions this is
multiplication by that character; it preserves cuspidality,
and twisting is an equivalence with its own inverse.
Its restriction to the connected real group is unchanged, so
its holomorphic weight is \(2e+1\).  The central character is
unchanged because \(\delta^2=1\).  At \(p\) the entire local
representation is unchanged.  Hence the associated newform
\(f_0^\delta\) still has underlying level prime to \(p\), the
same ordinary stabilization, and nebentype \(\eta_K\).
The equality of unramified Frobenius polynomials gives
\(V(f_0^\delta)=V(f_0)\otimes\delta\), by the semisimple
Galois representation dictionary and Chebotarev.
Put
\[
 V^\delta=V\otimes\delta,\qquad
 \gamma^\delta=\gamma\delta.
\]
It follows that
\[
 \det V^\delta=\eta_K,\qquad
 \overline{V^\delta}^{\rm ss}
   =(\eta_K\overline\gamma^\delta)
      \oplus(\overline\gamma^\delta)^{-1},\qquad
 (\overline\gamma^\delta)^2=\overline\gamma^2.      \tag{G81.6}
\]
Scalar twisting preserves irreducibility on \(G_K\).

Now choose the host of AC-G77 for this twisted input, putting
\(r_0\) into its prescribed splitting set.
Thus \(r_1\) is fresh for all the twisted input, \(K'\) splits
at \(p\) and at every prescribed old ramified prime, and
\[
 F=KK',\qquad F^+=\mathbf Q(\sqrt{3r_1}),\qquad F_0=K'
\]
satisfy all the source host and parity conditions.
The field \(F\) is unramified at \(r_0\): both \(K\) and \(K'\)
are unramified there.  Choose an element \(h_0\) of absolute
inertia at \(r_0\) for which \(\delta(h_0)=-1\).
This element belongs to \(G_F\).  The untwisted \(V\), \(\chi\)
and \(\epsilon\) are unramified at \(r_0\), and therefore
\[
       V^\delta(h_0)=-I_2,\qquad
       \overline\epsilon(h_0)=1,\qquad
       \chi(h_0)=1.                               \tag{G81.7}
\]
This is the literal scalar-witness premise of AC-G78; indeed it
holds before reduction.  It is not obtained by killing the
twist on the host.

For clarity, irreducibility of \(V^\delta\) persists on \(G_F\)
as required in the packet construction.  If the restriction
of the irreducible two-dimensional \(G_K\)-representation
were reducible under the quadratic extension \(F/K\), it would
admit the quadratic self-twist of \(F/K\).
That quadratic character ramifies at the fresh \(r_1\), whereas
\(V^\delta\) is unramified there, a contradiction.
This is the same fresh-prime argument proved in AC-G77.
Likewise (G81.1) persists on \(G_F\): otherwise the nontrivial
character \(\overline\gamma^2|_{G_K}\) would equal the quadratic
character of \(F/K\), impossible at \(r_1\).

\emph{Applying the actual ordinary construction to the new input.}
The proof of AC-G78 uses the displayed ordinary modular input,
determinant and residual shape, the host conditions, and the
literal sign alternative (G78.4).
It does not require \(\psi(-1)\) as an additional deformation
hypothesis.  Although \(\psi\delta\) has the opposite value at
complex conjugation, (G81.6) has the required global residual
shape, the rank-two determinant and oddness are unchanged,
and the ordinary parameter at \(p\) is identical to the old
one.  The endoscopic packet in Thorne's Theorem~4.5 uses the
same two-dimensional conjugate-self-dual parameter after
restriction and the same repeated parallel infinity types.
Consequently the proof applies to \(f_0^\delta,\tau\).
Here \(V^\delta\) is used solely as an ordinary automorphic
input; we are not constructing an old \(\mathrm{Dat}\) for it.

The genericity witness in AC-G78 is also retained.
The extra quadratic character is a rational character; one
still varies the anticyclotomic cubic value while fixing all
rational character values.  The quadratic host cannot kill a
cubic character, and its fresh ramification makes the finite
character conditions extend just as in AC-G77.
Equation (G81.7) supplies the alternative when
\(\tau(-1)(-1)^e=1\).  When the sign is negative, the other
alternative of AC-G78 already applies.

We spell out why the odd-weight application statements of
AC-G78-type and AC-G79 introduce no hidden use of odd \(e\).
In their proofs oddness supplies the sign witness and, in the
inherited application, guarantees residual separation.
We now have the first from (G81.7) and the second from (G81.1).
The free actual ordinary module, dense half-weight branch,
ramified-curve root construction, finite Hecke image and
arbitrary-base-change Cayley--Hamilton step are unchanged.
The same source-normalized exact full-chain fiber is obtained
at a rational \(q\), split completely in \(F\), with
\[
 q\equiv-1\pmod p,\qquad
 \chi(\operatorname{Frob}_q)=1.
\]
The four elliptic tests of AC-G75 therefore give an actual
ordinary Steinberg constituent at the original symmetric
classical weight.

At full \(q\)-level the integral ordinary residual-Hecke
lattice is projective over \(\mathcal O[\mathrm{GL}_3(k_q)]\).
The source residual genericity witness makes every
characteristic-zero split constituent generic.
A generic spherical constituent has finite restriction
\(\operatorname{Ind}_{B(k_q)}^{\mathrm{GL}_3(k_q)}1\);
its multiplicities of the finite Steinberg and trivial
representations are both one.  A nonspherical constituent
has trivial multiplicity zero and Steinberg multiplicity
nonnegative.  The actual smooth Steinberg contributes
strictly positively.  The projective Brauer-character
identity of AC-G72-character therefore gives an actual
finite type
\[
 \tau_\theta=
 \operatorname{Ind}_{P_{2,1}(k_q)}^{\mathrm{GL}_3(k_q)}
       (C_\theta^{(2)}\otimes1),
\]
where \(\theta\) is a regular character of order \(p\) of
\(k_{q^2}^{\times}\), trivial on \(k_q^\times\).
This repeats the proof of AC-G78-type with precisely its
separation and sign inputs now verified.

The local parameter has an irreducible tame two-dimensional
summand and an unramified line, with \(N=0\).
Its residual unitary Frobenius roots are
\(\{-1,1,-1\}\).  The two-dimensional summand has roots
\(a,-a\), so the line reduces to \(-1\).
The residual global characters are
\(\overline\chi^{\,j}\overline\gamma^\delta,1,
 \overline\chi^{\,j}(\overline\gamma^\delta)^{-1}\).
They are distinct.  Because of (G81.1), only the middle
character is individually conjugate self-dual.
A global \(2+1\) decomposition would therefore have a
one-dimensional conjugate-self-dual constituent reducing to
that middle character, contradicting its local Frobenius
value \(-1\).  A decomposition into three lines cannot
contain the local irreducible two-dimensional summand.
These exhaust the nontrivial dimension partitions.
Hence the normalized rank-three parameter \(R_F\) is
absolutely irreducible.

Finally impose the full-unit coefficient-prime projector of
AC-G79 throughout the actual construction.
Since \(\delta|_{G_{\mathbf Q_p}}=1\), its character is exactly
the same \(\delta_e\) as for the original \(V\).
The full-torus projector, actual free specialization, four
tests and projective \(q\)-type extraction commute exactly
as in that proof.  The ordered Newton slopes are
\((-fe,0,fe)\), and the finite inertia characters are
\((\chi^j\tau,1,\chi^j\tau)\).
The middle character differs from the endpoints; equal
endpoint characters have slope gap \(2ef\), whereas a
nonzero monodromy arrow would require gap \(f\).
Thus \(N=0\), also for \(e=1\).
The parallel weak-admissibility argument of AC-G79 gives
the precise flag (G81.4) over \(F\).

\emph{Killing the auxiliary twist only after irreducibility.}
Set \(L=FD\).  Fresh ramification at \(r_0\) shows \(D\not\subset F\),
so \(L/F\) is quadratic.
All three quadratic generators are imaginary, and complex
conjugation is their simultaneous sign change.
Consequently \(L\) is CM with \([L^+:\mathbf Q]=4\).
Since \(F\cap L^+=F^+\), one has \(L=FL^+\).
It follows that \(L/L^+\) is the base change of \(F/F^+\)
and is unramified at every finite place.
The imaginary quadratic subfield \(K'\) is not contained in
\(L^+\), so \(L=L^+K'\).
Since \(D\) splits at \(p\), every completion of \(L\) above a
given \(F\)-place over \(p\) equals that completion of \(F\).

The restriction \(R_F|_{G_L}\) is absolutely irreducible.
Here is a complete index-two argument, valid for the
potentially infinite Galois image.
If \(H\triangleleft G\) has index two and \(W\) is irreducible
finite-dimensional over an algebraically closed field of
characteristic zero, choose an irreducible \(H\)-submodule.
The sum of its two \(G/H\)-translates is a nonzero
\(G\)-submodule, hence all of \(W\); it is semisimple as an
\(H\)-module.  Frobenius reciprocity and the permutation
representation of \(G/H\) give
\[
 \operatorname{End}_{H}(W)
   \simeq
 \operatorname{Hom}_{G}(W,W\oplus W\otimes\delta_{L/F}).
                                                        \tag{G81.8}
\]
For \(\dim W=3\), a self-twist
\(W\simeq W\otimes\delta_{L/F}\) would imply
\(\delta_{L/F}^{\,3}=1\) by determinants, contradicting its
order two.  Schur's lemma therefore makes (G81.8)
one-dimensional.  Semisimplicity then forces the restriction
to be irreducible.  Apply this to \(W=R_F\).

The strong base change \(\Pi_F\) of the rank-three unitary
constituent is RACSDC.  More explicitly, the source residual
genericity excludes noncuspidal discrete (Speh) blocks by
Thorne's paragraph following Corollary~4.3, and irreducibility
of \(R_F\) excludes multiple cuspidal summands in Theorem~4.2.
Thus the strong base change is cuspidal.
Clozel--Thorne, \emph{Level raising and symmetric power
functoriality I}, Theorem~2.7(1), printed p.~6, supplies the
cuspidal soluble base change \(\Pi_L\).
Its Galois parameter is the restriction of the old one.
The printed proof reduces cuspidality to cyclic prime
degree and the absence of a self-twist; the required
absence was just proved.  Arthur--Clozel's base-change
existence theorem remains an imported input.

The representation \(\Pi_F\) is unramified at every place
inert over \(F^+\).  Ramified places were relative split
places, including the auxiliary \(q\) and the coefficient
prime.  Base change preserves unramified local parameters,
and a relative split place remains split after base change
from \(F^+\) to \(L^+\).  Hence every ramified place of
\(\Pi_L\) is split over \(L^+\).
Thorne's Theorem~4.4, printed p.~21, therefore descends
\(\Pi_L\) to an automorphic representation \(\sigma_L\) of
the definite unitary group for \(L/L^+\).
Its split local components equal those of \(\Pi_L\), and its
infinite weight is the unchanged parallel symmetric weight.
In particular the coefficient-prime components and flag
are unchanged, since the completions there are unchanged.
This proves ordinarity and (G81.4), without an inference
from residual ordinarity alone.

On \(G_L\) the character \(\delta\) is trivial.
Thus the residual parameter becomes exactly (G81.3).
Residual separation also persists on \(G_L\): if
\(\overline\gamma^2|_{G_F}\) became trivial there, it would
be the nontrivial quadratic character of \(L/F\).
The former is unramified at the fresh \(r_0\), whereas the
latter is ramified there, a contradiction.
The cubic character retains order three under extensions
of degree a power of two.
At \(q\) finite inertia and \(N=0\) persist on restriction.
The irreducible local two-dimensional summand may split
after \(L/F\); we do not use its irreducibility again.
Global irreducibility was established before base change
and preserved by (G81.8).

All arguments are qualitative.  The additional quadratic
host can introduce class-group sectors and change the
numerical constants of a deformation frame.  Although
restriction followed by corestriction multiplies a
\(269\)-primary cohomology class by a unit, that fact does
not verify any of the numerical frame inequalities or
per-curve automorphic support conditions.
The restrictions in the statement follow.
\end{proof}

\begin{remark}[Scope relative to the existing gates]
This proves a qualitative actual rank-three seed with the
specified coefficient-prime flag on the separated locus.
The original odd-\(e\) theorem is unchanged; for even \(e\)
the new feature is that a fresh odd quadratic character
supplies a literal scalar witness while it is still
nontrivial on the host.
It is killed only after global irreducibility has been
proved.  Any use of a further AC-G73 cover must separately
verify that cover's hypotheses.  Neither that cover nor
the quadratic base change is a substitute for the actual
\(J\)-containment required by AC-G61.
\end{remark}
